\documentclass[10pt]{article}
\usepackage[T1]{fontenc}
\usepackage[utf8]{inputenc}
\usepackage[french]{babel}
\usepackage{amsmath,amssymb,mathrsfs,mathpazo}
\usepackage{geometry}
\geometry{a4paper,margin=2cm}
\usepackage{graphicx,array,tabularx,tabularray,colortbl,multirow,multicol,booktabs,enumitem,pifont,ulem,fancybox,diagbox,cancel}
\usepackage{pstricks,pstricks-add,pst-plot,pst-tree,pst-node,pst-text,pst-func,pst-math,pst-bspline,pst-3dplot,pst-eucl,pst-solides3d}
\let\pscalculate\undefined
\usepackage{pst-bezier}
\usepackage{tikz,pgfplots,tkz-tab}
\usetikzlibrary{arrows,arrows.meta,patterns,calc,positioning,trees,decorations.markings,shapes.geometric,shapes.misc,shapes.arrows,intersections,decorations.pathmorphing,decorations.pathreplacing,backgrounds,shadows,babel}
\pgfplotsset{compat=1.18}
\usepgfplotslibrary{fillbetween}
\usepackage[np]{numprint}
\usepackage{listings,makecell,pas-tableur,twoopt}
\let\barre\undefined
\usepackage[pstricks]{tablvar}
\definecolor{darkgreen}{RGB}{0,100,0}
\providecommand{\red}{\color{red}}
\providecommand{\blue}{\color{blue}}
\providecommand{\green}{\color{green}}
\providecommand{\black}{\color{black}}
\providecommand{\white}{\color{white}}
\providecommand{\gray}{\color{gray}}
\providecommand{\cyan}{\color{cyan}}
\providecommand{\magenta}{\color{magenta}}
\providecommand{\yellow}{\color{yellow}}
\providecommand{\violet}{\color{violet}}
\providecommand{\orange}{\color{orange}}
\providecommand{\euro}{\texteuro}
\providecommand{\EUR}{\texteuro}
\providecommand{\ieme}{\textsuperscript{e}}
\providecommand{\degres}{\ensuremath{^\circ}}
\providecommand{\degree}{\ensuremath{^\circ}}
\providecommand{\vv}[1]{\overrightarrow{#1}}
\providecommand{\Coord}[1]{#1}
\providecommand{\decosix}{}
\providecommand{\starredbullet}{\ensuremath{\bullet}}
\pagestyle{empty}
\setlength{\parindent}{0pt}
\usepackage{adjustbox}\DeclareUnicodeCharacter{2236}{:}\DeclareUnicodeCharacter{2190}{\ensuremath{\leftarrow}}
\providecommand{\R}{\mathbb{R}}
\providecommand{\N}{\mathbb{N}}
\providecommand{\Z}{\mathbb{Z}}
\providecommand{\Q}{\mathbb{Q}}
\providecommand{\D}{\mathbb{D}}
\providecommand{\C}{\mathbb{C}}
\providecommand{\e}{{\mathrm{e}}}
\providecommand{\vect}[1]{\overrightarrow{#1}}
\providecommand{\vectt}[1]{\overrightarrow{\mathrm{#1}}}
\providecommand{\vvt}[1]{\overrightarrow{\mathrm{#1}}}
\providecommand{\np}[2][]{{#2~#1}}
\providecommand{\vv}[1]{{\overrightarrow{#1}}}
\renewcommand{\d}{{\mathrm{d}}}
\providecommand{\degres}{{^\circ}}
\providecommand{\degre}{{^\circ}}
\providecommand{\blue}{\color{blue}}
\providecommand{\red}{\color{red}}
\providecommand{\enskip}{\quad}
\providecommand{\llbracket}{\lbrack\!\lbrack}
\providecommand{\rrbracket}{\rbrack\!\rbrack}
\providecommand{\euro}{{\text{€}}}
\providecommand{\no}{{\text{n°}}}
\providecommand{\interval}[2]{{[#1;#2]}}
\providecommand{\intervFO}[2]{{[#1;#2[}}
\providecommand{\dotfill}{\ldots}

\providecommand{\ds}{\displaystyle}
\providecommand{\textding}[1]{#1}
\providecommand{\Oij}{\(\mathrm{O};\vec{\imath},\vec{\jmath}\)}
\providecommand{\Oijk}{\(\mathrm{O};\vec{\imath},\vec{\jmath},\vec{k}\)}
\providecommand{\Ouv}{\(\mathrm{O};\vec{u},\vec{v}\)}
\providecommand{\intervFF}[2]{\ensuremath{[#1;#2]}}
\renewcommand{\euro}{\ifmmode\text{\texteuro}\else\texteuro\fi}
\DeclareUnicodeCharacter{00B0}{\ifmmode^{\circ}\else\textdegree\fi}
\DeclareUnicodeCharacter{20AC}{\ifmmode\text{\texteuro}\else\texteuro\fi}
% Énoncés : Éducation nationale. Transcription et corrigés : APMEP, Denis Vergès et contributeurs. Découpage et classement : CoopMaths, Sébastien Lozano, Éric Elter et contributeurs. Adaptation de la présentation : Maths Explorer.
% CC BY-SA : https://coopmaths.fr/www/about/
% https://coopmaths.fr/alea/static/bac/2023/tex/bac_2023_03_sujet1_etrangers_groupe1_2.tex
\begin{document}
\noindent\textit{APMEP / CoopMaths. CC BY-SA. Présentation adaptée par Maths Explorer.}\par\medskip
\medskip

\textbf{Les deux parties de cet exercice sont indépendantes.}

\medskip

Dans une grande ville française, des trottinettes électriques sont mises à disposition des usagers. Une entreprise, chargée de l'entretien du parc de trottinettes, contrôle leur état chaque lundi.

\medskip

\textbf{Partie A}

\medskip

On estime que :
\begin{itemize}
\item lorsqu'une trottinette est en bon état un lundi, la probabilité qu'elle soit encore en bon état le lundi suivant est $0,9$ ;
\item lorsqu'une trottinette est en mauvais état un lundi, la probabilité qu'elle soit en bon état le lundi suivant est $0,4$.
\end{itemize}

\smallskip

On s'intéresse à l'état d'une trottinette lors des phases de contrôle.

Soit $n$ un entier naturel.
On note $B_n$ l'évènement \og la trottinette est en bon état $n$ semaines après sa mise en service\fg{} et $p_n$ la probabilité de $B_n$.

Lors de sa mise en service, la trottinette est en bon état. 
On a donc $p_0 = 1$.

\medskip

\begin{enumerate}
\item Donner $p_1$ et montrer que $p_2 = 0,85$.

\emph{On pourra s'appuyer sur un arbre pondéré.}
\item Recopier et compléter l'arbre pondéré ci-dessous :

\begin{center}
\pstree[treemode=R,nodesepA=0pt,nodesepB=2.5pt,treesep = 1cm,levelsep=2.5cm]{\TR{}}
{\pstree{\TR{$B_n$~}\taput{$p_n$}}
{\TR{$B_{n+1}$}\taput{\ldots}
\TR{$\overline{B_{n+1}}$}\tbput{\ldots}
}
\pstree{\TR{$\overline{B_n}$~}\tbput{\ldots}}
{\TR{$B_{n+1}$}\taput{\ldots}
\TR{$\overline{B_{n+1}}$}\tbput{\ldots}
}}

\end{center}

\item En déduire que, pour tout entier naturel $n$,\: $p_{n+1} 
= 0,5p_n + 0,4$.
\item 
	\begin{enumerate}
		\item Démontrer par récurrence que pour tout entier naturel $n$,\: $p_n \geqslant 0,8$.
		\item À partir de ce résultat, quelle communication l'entreprise peut-elle envisager pour valoriser la fiabilité du parc ?
	\end{enumerate}	
\item 
	\begin{enumerate}
		\item On considère la suite $\left(u_n\right)$ définie pour tout entier naturel $n$ par $u_n = p_n - 0,8$. 
		
Montrer que $\left(u_n\right)$ est une suite géométrique dont on donnera le premier terme
et la raison.
		\item En déduire l'expression de $u_n$ puis de $p_n$ en fonction de $n$.
		\item En déduire la limite de la suite $\left(p_n\right)$.
	\end{enumerate}
\end{enumerate}

\medskip

\textbf{Partie B}

\medskip

Dans cette partie, on modélise la situation de la façon suivante :
\begin{itemize}
\item l'état d'une trottinette est indépendant de celui des autres ;
\item la probabilité qu'une trottinette soit en bon état est égale à $0,8$.
\end{itemize}
\smallskip

On note $X$ la variable aléatoire qui, à un lot de $15$ trottinettes, associe le nombre de trottinettes en bon état.

Le nombre de trottinettes du parc étant très important, le prélèvement de 15 trottinettes peut être assimilé à un tirage avec remise.

\medskip

\begin{enumerate}
\item Justifier que $X$ suit une loi binomiale et préciser les paramètres de cette loi.
\item Calculer la probabilité que les $15$ trottinettes soient en bon état.
\item Calculer la probabilité qu'au moins $10$ trottinettes soient en bon état dans un
lot de $15$.
\item On admet que $E(X) = 12$. Interpréter le résultat.
\end{enumerate}

\end{document}
